\documentclass[11pt]{article}

\usepackage[margin=1in]{geometry}
\usepackage{booktabs}
\usepackage{hyperref}
\usepackage{xcolor}
\usepackage{parskip}
\usepackage{cablp-macros}

\hypersetup{colorlinks=true, linkcolor=blue, citecolor=blue, urlcolor=blue}

\title{\textbf{Pseudo-1D Transport Modeling in the LAPD}}
\author{LAPD Group}
\date{\today}

\begin{document}
\maketitle

\begin{abstract}
This manuscript describes the reduced transport model used to evolve plasma
and neutral dynamics in LAPD discharge simulations.
\end{abstract}

\section{Model Equations}

For each cell $i$, the evolved state is
\begin{equation}
  \bm{y}_i = \paren{\cell{\ne}{i}, \cell{\nn}{i}, \cell{\Te}{i},
  \cell{\Ti}{i}, \cell{\vplasma}{i}}.
\end{equation}

The density equations can be written compactly as
\begin{align}
  \deriv{\cell{\ne}{i}}{t}
    &= \cell{\Sionbulk}{i}
     + \cell{\Sionbeam}{i}
     - \cell{\Srecrad}{i}
     - \cell{\Srectb}{i}
     + \cell{\Fluxne}{i}, \\
  \deriv{\cell{\nn}{i}}{t}
    &= \cell{\Fluxnn}{i}
     - \paren{\cell{\Sionbulk}{i} + \cell{\Sionbeam}{i}
       - \cell{\Srecrad}{i} - \cell{\Srectb}{i}}
       \paren{\frac{\cell{\Rp}{i}}{\cell{\Rm}{i}}}^2
     + \cell{\Sgp}{i}
     - \cell{\Spump}{i}\cell{\nn}{i}.
\end{align}

The electron and ion temperature equations are
\begin{align}
  \deriv{\cell{\Te}{i}}{t}
    &= \cell{\Qeb}{i}
     - \cell{\Qie}{i}
     - \cell{\Qei}{i}
     - \cell{\Qen}{i}
     + \cell{\Qepara}{i}
     - \cell{\Qeperp}{i}
     + \cell{\Qdivv}{e,i}
     + \cell{\Qconv}{e,i}, \\
  \deriv{\cell{\Ti}{i}}{t}
    &= \cell{\Qie}{i}
     + \cell{\Qipara}{i}
     - \cell{\Qiperp}{i}
     - \cell{\Qcx}{i}
     + \cell{\Qdivv}{i,i}
     + \cell{\Qib}{i}
     + \cell{\Qconv}{i,i}.
\end{align}

Bulk electron-impact ionization is represented by
\begin{equation}
  \cell{\Sionbulk}{i}
    = b_{\mathrm{ioniz}}\,\cell{\ne}{i}\cell{\nn}{i}
      \Kion\paren{\cell{\Te}{i}; I}.
\end{equation}

\section{Macro Examples}

The reusable notation lives in \texttt{cablp-macros.sty}.  Typical examples are:
\begin{align}
  \Qcx &= \fE b_{Qcx}\,\nn\,\sgv_{\mathrm{cx}}(\Ti)\paren{\Ti-\Tn}, \\
  \kepara &= 3.16\,\taue \vthe^2, &
  \Qepara^{\mathrm{loss}}
    &= \fE b_{\mathrm{epara}}\frac{\kepara\Te}{\Lp\Lhf}.
\end{align}

\end{document}
